% FILE: 02_lineage_and_context.tex
% Project: otel-weaver
% Role: telemetry-to-process-evidence-ingestion-layer

\section{Lineage and Context}
\label{sec:otel-weaver-lineage}

\subsection{2016 Language-Model Lesson}
\label{sec:otel-weaver-lineage-2016}

The foundational intellectual event in this research program's lineage was a 2016 experiment
demonstrating that local text imitation does not conserve consequence. A language model trained
to reproduce the surface form of a process description---its vocabulary, sentence structure,
and topical coverage---failed to preserve the causal dependencies and lifecycle constraints
embedded in the original. The model could produce output that \emph{resembled} a process
specification without \emph{representing} the process's legal and temporal obligations.

This lesson maps directly onto the \texttt{otel-weaver} problem. An OpenTelemetry span resembles
a process event: it has a name, a timestamp, a duration, and contextual attributes. A Weaver
schema validation pass resembles a conformance verdict: it checks structure against a declared
template. But resemblance is not identity. The 2016 lesson established that surface-form fidelity
is insufficient for consequence-conserving representation. The \texttt{otel-weaver} doctrine layer
operationalizes this lesson at the ingestion stratum: a span that passes schema validation has
demonstrated structural fidelity, not process fidelity. The Feedstock Theorem
($\text{SchemaValidation}(\text{Telemetry}) \neq \text{ProcessCompliance}(\text{ProcessState})$)
is the type-system encoding of the 2016 finding.

The implication for the \texttt{otel-weaver} architecture is that no amount of schema richness
in the \texttt{process.pi} semantic convention registry can substitute for a formal admission
gate. The registry makes feedstock more \emph{legible}; it does not make feedstock
\emph{evidence}. That promotion requires a witness seal---the \texttt{ActivityWitness} struct
binding the admitted payload to a lattice-bounded accountability structure.

\subsection{Enterprise Process Gap}
\label{sec:otel-weaver-lineage-enterprise}

The enterprise process gap---the structural failure of deployed software systems to produce
auditable evidence of the processes they claim to execute---manifests acutely in organizations
that have adopted OpenTelemetry. These organizations possess rich observability infrastructure:
distributed traces, metric pipelines, structured log aggregation, and Weaver-validated schema
registries. Yet when asked to demonstrate that a specific business process executed in conformance
with its declared model, they cannot produce a replay-grade event log. They can produce dashboards.
They can produce alert histories. They cannot produce a process conformance receipt.

The \texttt{otel-weaver} project names this gap at the ingestion layer. The gap exists precisely
at the transition from OTel span to OCEL event. The \texttt{OtelToOcelLossPolicy} artifact
specifies three categories of loss that are mandatory to declare: \texttt{span\_duration\_truncation}
(temporal precision lost when mapping span timing to OCEL timestamps),
\texttt{resource\_scope\_collapse} (hierarchical resource context flattened into flat OCEL
attributes), and \texttt{causality\_to\_correlation\_degradation} (causal span relationships
degraded to correlational case identifiers). Each of these losses reduces the downstream
process-intelligence layer's ability to reconstruct lawful object lifecycles.

The enterprise process gap is therefore not a gap in instrumentation density---most organizations
have abundant telemetry---but a gap in \emph{evidence-grade projection}: the formal, loss-declared,
witness-sealed transformation from raw observability signal to process-intelligence feedstock.
\texttt{otel-weaver} is the research foundry's answer to this specific gap.

\subsection{Relationship to the Chatman Equation}
\label{sec:otel-weaver-lineage-chatman-equation}

The Chatman Equation $A = \mu(O^*)$ asserts that every asserted activity $A$ is the manufactured
output $\mu$ of an authoritative object state $O^*$. This equation governs the entire foundry's
claim-making architecture: no activity may be asserted without tracing to an authoritative object.
The \texttt{otel-weaver} ingestion layer is the mechanism by which raw telemetry is either
\emph{promoted} to a candidate $O^*$ or \emph{refused} before it can contaminate the equation's
input domain.

Concretely, an OTel span arrives at the ingestion layer as untyped feedstock. The
\texttt{process.pi} registry specifies what a span must contain to be a candidate activity
witness: a valid \texttt{instance\_id}, a declared \texttt{activity.name}, a \texttt{lifecycle}
state, token state fields (\texttt{before}/\texttt{after}), a \texttt{witness.id}, and a
\texttt{witness.hash} in BLAKE3 64-character format. If any required field is absent or malformed,
the \texttt{Admission<T,\,W>} gate emits a \texttt{Refusal} with a \texttt{RefusalCode} and
\texttt{violated\_rule} reference. The refused span never enters the equation's domain.

If the span passes, it is wrapped in an \texttt{Admission<T,\,W>} struct where \texttt{W} is an
\texttt{ActivityWitness} lattice element. This struct is the typed form of a candidate $O^*$. Only
at this point may downstream components---\texttt{wasm4pm}, \texttt{wasm4pm-compat}, the M\&A
claim layer---treat the payload as a process-evidence input. The Chatman Equation's integrity
depends on the ingestion gate holding this boundary without exception.

\subsection{Relationship to Receipts and Replay}
\label{sec:otel-weaver-lineage-receipts}

The foundry's receipt architecture requires that every claim be cryptographically anchored to an
immutable artifact. The \texttt{otel-weaver} project contributes to this architecture at two
levels. First, the \texttt{Refusal} type requires a \texttt{blake3\_digest} of the refused
feedstock payload, creating an auditable record of what was rejected and why. Second, the
\texttt{ggen}-manufactured manifest chain produces BLAKE3-chained receipt files at
\texttt{receipts/otel-weaver-source-receipt.json} and \texttt{receipts/pi-telemetry-bridge-receipt.json},
referenced in the project's two manifest files.

The Replay Test---a ground truth criterion requiring independent third-party replayability without
visualization---is the ultimate arbiter of whether the ingestion layer has succeeded. A process
event log derived through \texttt{otel-weaver}'s ingestion pipeline must support replay: a third
party, given the OCEL log and the declared process model, must be able to reproduce the conformance
verdict without access to dashboards, alerts, or observability infrastructure. This criterion is
what distinguishes a receipt from a reading. The \texttt{otel-weaver} ingestion layer is designed
so that every admitted payload carries enough information for replay, and every loss is declared
in a \texttt{LossReport} so that the replay analyst knows what was irreversibly discarded.
